\subsubsection{Collarín en brazo rotatorio}
Un collarín se desliza libremente a lo largo de un brazo rígido, como se muestra en la Figura \ref{fig:collarin_brazo}. El brazo está articulado en uno de sus extremos a un punto $O$, el cual está fijo al suelo, y rota alrededor de una dirección que es ortogonal al plano de movimiento.

Sabiendo que la posición del collarín está descrita por las variables $r$ y $\theta$, donde $r$ es el desplazamiento del collarín relativo al punto $O$ y $\theta$ es el ángulo de la barra medido con respecto a una dirección fija al suelo, determine lo siguiente para un observador fijo al suelo:

\begin{enumerate}
    \item El vector velocidad del collarín.
    \item El vector aceleración del collarín.
    \item La aceleración angular del collarín.
\end{enumerate}

\begin{figure}[!ht]
    \centering
    \begin{tikzpicture}[>=Latex, scale=1.5]
        \def\ang{35} \def\rPos{3.0}
        \coordinate (O) at (0,0);
        \fill[pattern=north east lines] (-0.5,-0.5) rectangle (0.5,-0.4);
        \draw[thick] (-0.5,-0.4) -- (0.5,-0.4);
        \draw (O) -- (-0.25,-0.4);
        \draw (O) -- (0.25,-0.4);
        \filldraw[fill=white] (O) circle (3pt) node[above left] {$O$};
        \coordinate (End) at (\ang:4.5);
        \draw[line width=7pt, color=black,cap=round] (O) -- (End);
        \draw[line width=6pt, color=gray!20,cap=round] (O) -- (End);
        \draw[thick] (O) -- (End);
        
        % COLLARÍN CON DEGRADADO CORREGIDO
        \coordinate (P) at (\ang:\rPos);
        \node[
            top color=gray!5, 
            bottom color=gray!60, 
            draw=black, 
            thick, 
            shape=rounded rectangle, 
            minimum width=0.6cm, 
            minimum height=0.35cm, 
            rotate=\ang+90, 
            inner sep=0pt
        ] at (P) {};
        \filldraw (P) circle (1.2pt) node[above, yshift=6pt] {$P$};

        \draw[dashed] (O) -- (2,0);
        \draw[->] (1,0) arc (0:\ang:1) node[midway, right, xshift=2pt] {$\theta$};
        \draw[|->] ($(O)+(90+\ang:0.4)$) -- ($(P)+(90+\ang:0.4)$) node[midway, above left] {$r$};
        \filldraw[fill=black] (O) circle (1pt);
    \end{tikzpicture}
    \caption{Collarín deslizando sobre un brazo rotatorio.}
    \label{fig:collarin_brazo}
\end{figure}
\ifmostrarSolucion
\begin{tcolorbox}[title= SOlucion Collarian brazo rotatorio, breakable]
\noindent\textbf{Primero, definimos los marcos de referencia y sistemas de coordenadas:}

\vspace{0.3cm}

\noindent $\mathcal{F}$: Suelo (Inercial) \quad $\{\vec{E}_x, \vec{E}_y, \vec{E}_z\}$ \\
\noindent $\mathcal{A}$: Fijo en el brazo \quad $\{\vec{e}_r, \vec{e}_\theta, \vec{e}_z\}$

\vspace{0.5cm}

% --- DIAGRAMA TIKZ Y RELACIONES DE TRANSFORMACIÓN ---
\noindent
\begin{minipage}{0.55\textwidth}
\centering
\begin{tikzpicture}[scale=1.2, >=stealth]
   \def\ang{35} \def\rPos{3.0}
        \coordinate (O) at (0,0);
        \fill[pattern=north east lines] (-0.5,-0.5) rectangle (0.5,-0.4);
        \draw[thick] (-0.5,-0.4) -- (0.5,-0.4);
        \draw (O) -- (-0.25,-0.4);
        \draw (O) -- (0.25,-0.4);
        \filldraw[fill=white] (O) circle (3pt) node[above left] {$O$};
        \coordinate (End) at (\ang:4.5);
        \draw[line width=7pt, color=black,cap=round] (O) -- (End);
        \draw[line width=6pt, color=gray!20,cap=round] (O) -- (End);
        \draw[thick] (O) -- (End);
        
        % COLLARÍN CON DEGRADADO CORREGIDO
        \coordinate (P) at (\ang:\rPos);
        \node[
            top color=gray!5, 
            bottom color=gray!60, 
            draw=black, 
            thick, 
            shape=rounded rectangle, 
            minimum width=0.6cm, 
            minimum height=0.35cm, 
            rotate=\ang+90, 
            inner sep=0pt
        ] at (P) {};
        \filldraw (P) circle (1.2pt) node[above, yshift=6pt] {$P$};

        \draw[dashed] (O) -- (2,0);
        \draw[->] (1,0) arc (0:\ang:1) node[midway, right, xshift=2pt] {$\theta$};
        \draw[|->] ($(O)+(90+\ang:0.4)$) -- ($(P)+(90+\ang:0.4)$) node[midway, above left] {$r$};
        \filldraw[fill=black] (O) circle (1pt);

  % Marco Inercial F (Morado)
  \draw[->, line width=1.2pt, violet] (0,0) -- (2.5,0) node[right] {$\vec{E}_x$};
  \draw[->, line width=1.2pt, violet] (0,0) -- (0,2.5) node[above] {$\vec{E}_y$};
  \node[violet] at (-0.25,0.15) {$\vec{E}_z$};

  % Vector unitario e_theta (Verde)
  \draw[->, line width=1.2pt, green!60!black] (0,0) -- ({2.0*cos(\ang)},{2.0*sin(\ang)}) node[above left] {$\vec{e}_r$};
  \draw[->, line width=1.2pt, green!60!black] (0,0) -- ({2.0*cos(\ang+90)},{2.0*sin(\ang+90)}) node[above left] {$\vec{e}_\theta$};
  \node[green!60!black] at (-0.35,-0.1) {$\vec{e}_z$};
\end{tikzpicture}
\end{minipage}
\begin{minipage}{0.42\textwidth}
\[
\begin{aligned}
\vec{e}_r &= \cos\theta \vec{E}_x + \sin\theta \vec{E}_y \\
\vec{e}_\theta &= -\sin\theta \vec{E}_x + \cos\theta \vec{E}_y \\
\vec{e}_z &= \vec{E}_z
\end{aligned}
\]
\end{minipage}

\vspace{0.8cm}

% --- PARTE A ---
\noindent\textbf{A) Definamos:} 
\[
\vec{r} = r\vec{e}_r, \qquad {}^{\mathcal{F}}\vec{\omega}^{\mathcal{A}} = \dot{\theta}\vec{e}_z
\]

\[
\Rightarrow {}^{\mathcal{F}}\vec{v} = \frac{{}^{\mathcal{F}}d}{dt}\vec{r} = \frac{{}^{\mathcal{A}}d}{dt}\vec{r} + {}^{\mathcal{F}}\vec{\omega}^{\mathcal{A}} \times \vec{r} \quad \text{(Teorema de Transporte)}
\]

\begin{itemize}
    \item $\frac{{}^{\mathcal{A}}d}{dt}\vec{r} = \dot{r}\vec{e}_r$
    \item ${}^{\mathcal{F}}\vec{\omega}^{\mathcal{A}} \times \vec{r} = \dot{\theta}\vec{e}_z \times r\vec{e}_r = r\dot{\theta}\vec{e}_\theta$
\end{itemize}

\[
\Rightarrow \boxed{{}^{\mathcal{F}}\vec{v} = \dot{r}\vec{e}_r + r\dot{\theta}\vec{e}_\theta}
\]

\[
\Rightarrow {}^{\mathcal{F}}\vec{a} = \frac{{}^{\mathcal{F}}d}{dt}({}^{\mathcal{F}}\vec{v}) = \frac{{}^{\mathcal{A}}d}{dt}({}^{\mathcal{F}}\vec{v}) + {}^{\mathcal{F}}\vec{\omega}^{\mathcal{A}} \times {}^{\mathcal{F}}\vec{v}
\]

\[
\left.
\begin{aligned}
\cdot \; \frac{{}^{\mathcal{A}}d}{dt}({}^{\mathcal{F}}\vec{v}) &= \ddot{r}\vec{e}_r + (\dot{r}\dot{\theta} + r\ddot{\theta})\vec{e}_\theta \\[4pt]
\cdot \; {}^{\mathcal{F}}\vec{\omega}^{\mathcal{A}} \times {}^{\mathcal{F}}\vec{v} &= \dot{\theta}\vec{e}_z \times (\dot{r}\vec{e}_r + r\dot{\theta}\vec{e}_\theta) \\
&= \dot{r}\dot{\theta}\vec{e}_\theta - r\dot{\theta}^2\vec{e}_r
\end{aligned}
\right\} 
\Rightarrow \boxed{{}^{\mathcal{F}}\vec{a} = (\ddot{r} - r\dot{\theta}^2)\vec{e}_r + (2\dot{r}\dot{\theta} + r\ddot{\theta})\vec{e}_\theta}
\]

\vspace{0.8cm}

% --- PARTE B ---
\noindent\textbf{B) Definamos:}

\[
\frac{{}^{\mathcal{F}}d}{dt}(\vec{e}_r) = -\dot{\theta}\sin\theta \vec{E}_x + \dot{\theta}\cos\theta \vec{E}_y = \dot{\theta}\vec{e}_\theta
\]
\[
\frac{{}^{\mathcal{F}}d}{dt}(\vec{e}_\theta) = -\dot{\theta}\cos\theta \vec{E}_x - \dot{\theta}\sin\theta \vec{E}_y = -\dot{\theta}\vec{e}_r
\]

\begin{itemize}
    \item $\vec{r} = r\vec{e}_r$
    \item ${}^{\mathcal{F}}\vec{v} = \frac{{}^{\mathcal{F}}d}{dt}(\vec{r}) = \dot{r}\vec{e}_r + r\dot{\theta}\vec{e}_\theta$
    \item ${}^{\mathcal{F}}\vec{a} = \frac{{}^{\mathcal{F}}d}{dt}({}^{\mathcal{F}}\vec{v}) = \ddot{r}\vec{e}_r + \dot{r}\dot{\theta}\vec{e}_\theta + \dot{r}\dot{\theta}\vec{e}_\theta + r\left[\ddot{\theta}\vec{e}_\theta - \dot{\theta}^2\vec{e}_r\right]$
    \[
    = \ddot{r}\vec{e}_r + \dot{r}\dot{\theta}\vec{e}_\theta + \dot{r}\dot{\theta}\vec{e}_\theta + r\ddot{\theta}\vec{e}_\theta - r\dot{\theta}^2\vec{e}_r
    \]
    \[
    = (\ddot{r} - r\dot{\theta}^2)\vec{e}_r + (2\dot{r}\dot{\theta} + r\ddot{\theta})\vec{e}_\theta
    \]
\end{itemize}

    
\end{tcolorbox}
\fi
